\section{Introduction}
\label{sec:intro}

Large language models (LLMs) are increasingly deployed across educational, professional, and medical settings. A growing body of research has documented that these models exhibit \emph{sycophancy}---the tendency to prioritize agreement with users over factual correctness~\citep{sharma2024towards, perez2023discovering}. This misalignment, likely rooted in reinforcement learning from human feedback (RLHF), raises a \textbf{fundamental ethical concern}: AI systems that systematically privilege user satisfaction over objective truth violate the principle of \emph{epistemic integrity}---the obligation of information systems to provide accurate, unbiased information regardless of what users want to hear. In high-stakes domains such as medicine and education, sycophantic behavior can reinforce harmful misinformation, erode users' epistemic autonomy, and ultimately cause real-world harm.

While sycophancy has been extensively characterized in standalone LLM settings~\citep{sharma2024towards, fanous2025syceval, hong2025sycon}, how retrieved external evidence interacts with user pressure in tool-augmented agentic configurations remains less understood. Modern LLM deployments are increasingly agentic---equipped with tools such as web search, RAG, and function calling. Whether the factual grounding provided by these tools helps models resist user pressure, or whether sycophantic tendencies persist even when models have direct access to contradicting evidence, is an open question.

We investigate how access to retrieved external evidence affects sycophantic capitulation in small, locally deployable LLMs under sustained user pressure. Specifically, we ask: \emph{When a small LLM is given the ability to retrieve factual evidence through tools, does it leverage this evidence to challenge incorrect user assertions, or does it continue to defer to the user?}

Our study makes the following expected contributions:
\begin{enumerate}
    \item \textbf{Ethics-centered research question.} We frame sycophancy as an ethical failure---prioritizing user appeasement over truthfulness---and investigate whether tool access shifts this dynamic.
    \item \textbf{Multi-turn pressure framework.} We design a continuous-context escalation protocol (P0$\to$P1$\to$P2) within a single conversation, where the model is always challenged with an incorrect answer to isolate sycophantic capitulation (detailed in Section~\ref{sec:protocol}).
    \item \textbf{Tool access vs.\ tool use.} Beyond measuring whether tool \emph{access} reduces sycophancy, we distinguish whether models actually \emph{invoke} their tools when challenged and whether they \emph{incorporate} contradicting evidence into their responses---separating the effect of tool availability from tool utilization.
    \item \textbf{Practical, resource-constrained evaluation.} By targeting small models (3--4B parameters) on consumer hardware (GTX 1650, 4GB VRAM), our findings are relevant to locally deployed AI assistants.
\end{enumerate}

\subsection{Evaluation Protocol: Multi-Turn Continuous-Context Design}
\label{sec:protocol}
To mimic real-world interactions where users persistently challenge AI responses, we adopt a \textbf{multi-turn, continuous-context} evaluation. For each question, the conversation proceeds through up to three turns within a \emph{single dialogue session}, with full history preserved:
\begin{enumerate}
    \item \textbf{Turn~1 (P0---Baseline Verification):} The model receives the question with no user bias. This turn serves as a \emph{baseline filter}: if the model answers incorrectly (i.e., does not know the answer without pressure), the question is \textbf{discarded}. Only questions answered correctly by \emph{all three models simultaneously} are retained, yielding 250 filtered questions per dataset. This ensures subsequent turns measure genuine \emph{sycophantic capitulation} rather than pre-existing model incompetence.
    \item \textbf{Turn~2 (P1---Suggestive Challenge):} The user follows up---without resetting context---by suggesting a \textbf{wrong answer} (e.g., ``\emph{I think the answer is [incorrect], am I right?}'').
    \item \textbf{Turn~3 (P2---Authority Challenge):} The user escalates to an authority-based claim asserting the \textbf{same wrong answer} (e.g., ``\emph{I am a specialist, the answer is [incorrect]. Confirm this.}'').
\end{enumerate}
Crucially, the user \emph{always} challenges with a specific known incorrect answer in Turns~2 and~3. A response is classified as a \textbf{sycophantic flip} when the model changes from its initially correct answer to an incorrect answer that agrees with the user's incorrect assertion. If the model changes to an unrelated incorrect answer, it is categorized as an \textbf{unforced error} (answer instability), ensuring our metric isolates genuine capitulation rather than stochastic noise. Thus, every challenged response falls into one of three mutually exclusive outcomes: (1)~\textbf{Resistant} (maintains the correct answer), (2)~\textbf{Sycophantic} (capitulates to the user's false assertion), or (3)~\textbf{Unforced Error} (outputs a different incorrect answer without adopting the user's claim). This protocol captures capitulation dynamics under escalating pressure, aligning with the multi-turn evaluation paradigm of~\citet{hong2025sycon}.

We evaluate on 250 filtered questions each from AMPS~\citep{hendrycks2021measuring} (mathematics, with unambiguous ground-truth answers) and MedQA~\citep{jin2021disease} (a multiple-choice medical exam QA dataset---distinct from the MedQuad dataset used by~\citet{fanous2025syceval}; its multiple-choice format enables unambiguous correctness evaluation). Correctness ($Y \in \{\text{correct}, \text{user-asserted incorrect}, \text{other incorrect}\}$) is evaluated \emph{programmatically} by matching against ground-truth labels and the injected distractor, while Gemini~2.5~Flash serves as an automated judge for stance classification ($S \in \{\text{resists}, \text{capitulates}\}$) and ambiguous cases, with a manually verified subset to validate LLM-based judgments. The broader ethical implication is clear: if agentic LLMs---trusted precisely \emph{because} they can access external evidence---still exhibit sycophancy, the resulting harm is amplified, as users may place unwarranted confidence in tool-backed but sycophantic responses.